\label{sec:task}

% >>> Vicenc: rewritten, we don't mention the transformers yet nor f_\theta, just the supervision signals

We consider two tasks for learning an unknown lifted \strips model
$M=\langle\mathcal{P},A\rangle$ from traces drawn from problem instances
$P=\langle M,O_N,I,G\rangle$. In both cases, a trace contains an initial
state $s_0=I$ and a sequence $\tau=(a_0,\ldots,a_{n-1})$ of ground actions
obtained by instantiating the action schemas in $A$ with objects from $O_N$.
The two tasks differ in the supervision provided with these traces.

\subsection{Applicability Supervision}

In this setting, every action $a_i$ in the trace is labeled as applicable or
inapplicable given $s_0$ and the preceding actions. Action $a_0$ is applicable
when its preconditions hold in $s_0$, while $a_i$, for $i>0$, is applicable
when its preconditions hold in the state obtained from $s_0$ after processing
$a_0,\ldots,a_{i-1}$.

For a sequence $\tau_i=(a_0,\ldots,a_i)$, we define
$$
f_M(s_0,\tau_i)=
\begin{cases}
0, & \text{if $a_i$ is applicable after processing $a_0,\ldots,a_{i-1}$},\\
1, & \text{otherwise}.
\end{cases}
$$
The supervision therefore directly specifies the values of the applicability function \(f_M\) along the observed traces.

To define labels after an inapplicable action, we treat such actions as
no-ops: they leave the current state unchanged. Thus, applicable actions
update the state according to $M$, whereas inapplicable actions do not affect
subsequent applicability.
A trace is \emph{positive} if all its actions are applicable and
\emph{negative} otherwise. Effectively, each labeled sequence represents a
collection of traces: for every position $i$, removing the preceding
inapplicable actions yields an applicable prefix, and appending $a_i$ gives a
positive or negative trace according to its label. In our experiments, these
sequences are generated by random exploration while keeping applicable and
inapplicable examples balanced.

\subsection{Final-State Supervision}

This second setting is closer to world-model learning through state
prediction~\citep{world-models2,xie-etal-2025-making,spies2025transformers}.
Each trace contains an initial state $s_0$, an applicable action sequence
$\tau=(a_0,\ldots,a_{n-1})$, and the resulting final state $s_F$.
Intermediate states are not observed.

The supervision therefore specifies the state reached after executing the
complete action sequence, rather than the applicability of individual
actions. In particular, it provides the truth value in $s_F$ of each ground
predicate instance, while all actions in $\tau$ are known to be applicable.
The learning task is to infer a model whose dynamics map $s_0$ and $\tau$ to
$s_F$.
